% MỞ ĐẦU (LỜI MỞ ĐẦU - Chương 0)

\chapter*{MỞ ĐẦU} % Tên của chương
\addcontentsline{toc}{chapter}{MỞ ĐẦU} % Thêm tên chương vào mục lục

\label{Chapter0} % Để trích dẫn chương này ở chỗ nào đó trong bài, hãy sử dụng lệnh \ref{Chapter0} 

%----------------------------------------------------------------------------------------

Trong bối cảnh chuyển đổi số đang diễn ra mạnh mẽ cùng với sự phát triển vượt bậc của công nghệ thông tin, nhu cầu số hóa tài liệu để lưu trữ, quản lý và tìm kiếm ngày càng trở nên cấp thiết \cite{Reference1}. Tuy nhiên, trên thực tế, một lượng lớn tài liệu, đặc biệt là văn bản viết tay và văn bản tiếng Việt, vẫn chưa được số hóa hiệu quả. Các phương pháp nhập liệu thủ công không chỉ tốn nhiều thời gian, công sức mà còn dễ phát sinh sai sót, gây khó khăn trong việc khai thác và xử lý thông tin \cite{Reference2}.

Xuất phát từ thực tế đó, em lựa chọn đề tài “Nghiên cứu hiệu quả các kiến trúc học sâu trong bài toán nhận diện văn bản tiếng Việt”. Đề tài tập trung vào việc nhận diện chữ viết tay tiếng Việt và  khảo sát, xây dựng và so sánh hiệu quả của các mô hình học sâu tiêu biểu trong lĩnh vực nhận diện văn bản, nhằm tìm ra hướng tiếp cận phù hợp cho dữ liệu tiếng Việt – một ngôn ngữ có đặc thù về dấu và cấu trúc ký tự phức tạp \cite{Reference3}.

Trong báo cáo này, em trình bày quá trình nghiên cứu, thiết kế mô hình, huấn luyện và đánh giá các kiến trúc học sâu khác nhau áp dụng cho bài toán nhận diện văn bản chữ viết tay. Qua đó, đề tài hướng tới việc đưa ra những nhận xét, so sánh và đề xuất giải pháp nhằm nâng cao hiệu quả nhận diện, góp phần cải thiện chất lượng các hệ thống số hóa tài liệu tiếng Việt trong thực tế.

